% EXPANDED edition of the QIQT-H blueprint content.
% Same statements and \lean/\leanok/\uses wiring as content.tex, but the proofs
% are written as fuller, structure-mirroring arguments that follow the actual Lean
% proof (named helper lemmas, the real case splits and reductions). The machine
% checks remain the source of truth; this edition is for a reader who wants to see
% how each proof goes without opening the Lean files.
% Used by web-expanded.tex / print-expanded.tex. Do NOT add \begin{document} here.

\chapter{Overview (expanded edition)}

This is the \emph{expanded} rendering of the machine-checked core of \emph{QIQT-H}:
a single-world account of quantum measurement in which a unique outcome and the Born
statistics of repeated trials are \emph{derived} rather than postulated. It carries
the same numbered results as the concise edition, but each proof is spelled out to
mirror the underlying Lean argument.

\medskip
\noindent\textbf{Reading conventions.}
\begin{itemize}
  \item A green ``Lean'' tag means the statement is formalized and kernel-checked
        (no \texttt{sorry}); a green proof bubble means the proof is, too. The prose
        below names the actual lemmas the Lean proof uses, so it is a faithful map of
        the formal proof, not an independent re-derivation.
  \item Layered structure: \textbf{A} derives a single actual outcome from a
        finite-capacity record structure; \textbf{B} establishes Born as a conditional
        theorem; \textbf{C} builds a covariant typicality measure over histories;
        \textbf{Phase~B} supplies a normal state on $\BH$ and a free-field instance.
  \item Cited (not proved here) inputs --- e.g.\ that specific local quantum field
        algebras are type~III$_1$ --- are never tagged green and are flagged in text.
\end{itemize}


\chapter{Layer A: A Single Outcome Without Collapse}
\label{chap:layerA}

A measurement record is modelled by a finite \emph{record context} $C$: a type of
records $\Rec(C)$, a positive cost $\mathrm{cost}: \Rec(C)\to\mathbb{R}$ with each
cost exceeding half the capacity ($\mathrm{cost}(r) > \tfrac12\,\mathrm{capacity}$),
and a \emph{co-actual} configuration whose total active cost does not exceed the
capacity. A \emph{selection} $S$ additionally marks a nonempty active set. The
capacity bound forbids two incompatible records from being jointly active; the
selection forces at least one. Hence exactly one record is actual --- with no
collapse dynamics.

\begin{theorem}[Exactly one actual record]
  \label{thm:exactly-one-actual}
  \lean{QIQTH.CoreNoCollapse.exactly_one_actual}
  \leanok
  Let $C$ be a record context and $S$ a selection over $C$. Then the active set of
  $S$ has cardinality exactly one:
  \[
    \bigl|\, \mathrm{active}(S) \,\bigr| = 1 .
  \]
\end{theorem}

\begin{proof}
  \leanok
  Two bounds are combined.

  \emph{Upper bound} $|\mathrm{active}(S)| \le 1$ (lemma
  \texttt{coactual\_subsingleton}). Suppose for contradiction the cardinality were
  $\ge 2$. By \texttt{Finset.one\_lt\_card\_iff} there are two distinct active records
  $r_1 \ne r_2$, so $\{r_1,r_2\} \subseteq \mathrm{active}(S)$. Since costs are
  positive, summing over a superset only increases the total
  (\texttt{Finset.sum\_le\_sum\_of\_subset\_of\_nonneg}); together with
  \texttt{Finset.sum\_pair} this gives
  \[
    \mathrm{cost}(r_1) + \mathrm{cost}(r_2)
      \;\le\; \sum_{r \in \mathrm{active}(S)} \mathrm{cost}(r)
      \;\le\; \mathrm{capacity}.
  \]
  But each cost exceeds half the capacity ($\mathrm{cost\_gt\_half}$), so the
  left-hand side strictly exceeds the capacity --- a contradiction
  (closed by \texttt{linarith}).

  \emph{Lower bound} $1 \le |\mathrm{active}(S)|$: the selection's active set is
  nonempty, i.e.\ has positive cardinality (\texttt{S.selected.card\_pos}).

  Two natural-number inequalities $|\,\cdot\,| \le 1$ and $1 \le |\,\cdot\,|$ force
  equality (\texttt{omega}).
\end{proof}

\begin{corollary}[Single outcome, no collapse]
  \label{thm:single-outcome}
  \lean{QIQTH.CoreNoCollapse.qiqth_single_outcome_no_collapse}
  \leanok
  \uses{thm:exactly-one-actual}
  For every selection $S$ there is a \emph{unique} active record:
  \[
    \exists!\, r \in \Rec(C), \quad r \in \mathrm{active}(S).
  \]
  A definite outcome obtains without any collapse postulate.
\end{corollary}

\begin{proof}
  \leanok
  By Theorem~\ref{thm:exactly-one-actual} the active set has cardinality one, so by
  \texttt{Finset.card\_eq\_one} it equals a singleton $\{r\}$ for some $r$. Then
  membership in $\mathrm{active}(S)$ is membership in $\{r\}$
  (\texttt{Finset.mem\_singleton}): $r$ itself is active, and any active $y$ equals
  $r$. This is exactly unique existence.
\end{proof}


\chapter{Layer B: The Born Rule as a Conditional Theorem}
\label{chap:layerB}

Throughout, $\rho$ is a density matrix on $\mathbb{C}^d$ and $(E_k)_k$ a POVM;
$\Born(\rho,E)_k$ is the Born weight of outcome $k$, and for a length-$n$ outcome
sequence $\omega$, $w_{\rho,E}(\omega)$ is the quantum weight.

\begin{theorem}[Frequency concentration (quantum Chebyshev bound)]
  \label{thm:quantum-chebyshev}
  \lean{QIQTH.BornTypicalityQuantum.quantum_chebyshev_freq}
  \leanok
  Fix an outcome $k$, tolerance $\varepsilon>0$ and trial count $n>0$. With
  $p_k = \Born(\rho,E)_k$ and $\mathrm{count}_k(\omega)$ the number of occurrences of
  $k$ in $\omega$,
  \[
    \re \!\!\sum_{\substack{\omega \,:\, (n\varepsilon)^2 \,\le\,
        (\mathrm{count}_k(\omega) - n p_k)^2}}
    \! w_{\rho,E}(\omega)
    \;\le\; \frac{p_k\,(1-p_k)}{n\,\varepsilon^{2}} .
  \]
\end{theorem}

\begin{proof}
  \leanok
  The argument has two steps.

  \emph{Reduction to a real distribution.} For Hermitian $\rho$ and $E_k$, each
  quantum weight is real and equals the classical product weight of the Born
  distribution: $w_{\rho,E}(\omega) = w(\Born(\rho,E))(\omega)$
  (\texttt{quantumWeight\_eq\_w}). Summing this identity over the (finite) atypical set
  and taking real parts --- using $\re$ of a real-coerced sum
  (\texttt{Complex.ofReal\_sum}, \texttt{Complex.ofReal\_re}) --- rewrites the
  left-hand side as the real sum $\sum_\omega w(\Born(\rho,E))(\omega)$ over the same
  set.

  \emph{Classical Chebyshev.} The Born weights form a genuine probability
  distribution: nonnegative (\texttt{hp0}) and summing to one (\texttt{hp1}). The
  finite Chebyshev bound for i.i.d.\ trials
  (\texttt{BornTypicalityFinite.chebyshev\_freq}) applied to this distribution, with
  per-trial mean $p_k$ and variance $p_k(1-p_k)$, bounds the total weight of the
  deviation set $\{(n\varepsilon)^2 \le (\mathrm{count}_k - n p_k)^2\}$ by
  $p_k(1-p_k)/(n\varepsilon^2)$, which is the claim.
\end{proof}

\begin{theorem}[Uniqueness of the Born measure]
  \label{thm:born-unique}
  \lean{QIQTH.BornMeasureUniqueness.product_born_measure_unique}
  \leanok
  If $\mu$ is additive on events ($\mu(\varnothing)=0$ and
  $\mu(\mathrm{insert}\,a\,S)=\mu\{a\}+\mu(S)$ for $a\notin S$) with singleton weights
  $\mu\{\omega\} = \prod_t \Born(\rho,E)_{\omega_t}$, then
  \[
    \mu(S) \;=\; \re \,\Tr\!\bigl( \rho^{\otimes n}\, F_S \bigr)
    \qquad\text{for every event } S .
  \]
\end{theorem}

\begin{proof}
  \leanok
  First, $\mu$ coincides with the canonical \texttt{bornMeasure}. Both vanish on
  $\varnothing$ ($\mu$ by hypothesis, the other by \texttt{bornMeasure\_empty}); both
  are additive under \texttt{insert} ($\mu$ by hypothesis, the other by
  \texttt{bornMeasure\_insert}); and on singletons the hypothesis $\mu\{\omega\} =
  \prod_t \Born(\rho,E)_{\omega_t}$ matches \texttt{bornMeasure\_singleton}. A measure
  determined on singletons by additivity is unique
  (\texttt{measure\_unique\_of\_additive}), so $\mu = \texttt{bornMeasure}\,\rho\,E$
  on all finite events.

  Second, the canonical Born measure is the trace functional: for positive
  semidefinite (hence Hermitian) $\rho$ and $E_k$,
  $\texttt{bornMeasure}\,\rho\,E\,(S) = \re\,\Tr(\rho^{\otimes n} F_S)$
  (\texttt{bornMeasure\_eq\_trace}). Composing the two gives the stated identity.

  (Note: the product hypothesis silently bundles \emph{independence} across trials.
  This is unavoidable --- single-trial marginals alone do not force the product, since
  a maximally correlated measure shares them --- and is named explicitly, not buried,
  in the companion lemma \texttt{product\_born\_measure\_unique\_of\_independent\_trials}.)
\end{proof}


\chapter{Layer C: A Covariant Typicality Measure}
\label{chap:layerC}

Layer~C lifts the finite Born marginals to a single $\sigma$-additive probability
measure on the space of global selection histories $\prod_i \alpha_i$, consistent
(a projective limit), no-signaling, and state-agnostic.

\begin{definition}[No-signaling, bipartite]
  \label{def:no-signaling}
  A family of local statistics is \emph{no-signaling} if a local marginal never
  depends on the choice of a spacelike-remote measurement.
\end{definition}

\begin{theorem}[Bipartite no-signaling for arbitrary states]
  \label{thm:no-signaling}
  \lean{QIQTH.NoSignalingGeneral.bipartite_no_signaling}
  \leanok
  \uses{def:no-signaling}
  For any (possibly entangled) bipartite state $\rho$ on
  $\mathbb{C}^{d_1}\otimes\mathbb{C}^{d_2}$, local effect $E$, and remote POVM
  $(F_b)_b$ with $\sum_b F_b = 1$,
  \[
    \sum_b \Tr\!\bigl(\rho\,(E\otimes F_b)\bigr)
    \;=\; \Tr\!\bigl(\rho\,(E\otimes 1)\bigr) .
  \]
\end{theorem}

\begin{proof}
  \leanok
  Pull the finite sum inside the trace and factor it out of the product
  (\texttt{Matrix.trace\_sum}, \texttt{Finset.mul\_sum}):
  $\sum_b \Tr(\rho\,(E\otimes F_b)) = \Tr\bigl(\rho \sum_b (E\otimes F_b)\bigr)$.
  The Kronecker product is right-additive, so by induction on the index set
  $\sum_b (E\otimes F_b) = E \otimes \bigl(\sum_b F_b\bigr)$
  (\texttt{kronecker\_sum\_right}, proved by \texttt{Finset.induction} from
  \texttt{kronecker\_add}). Finally $\sum_b F_b = 1$ by hypothesis, leaving
  $\Tr(\rho\,(E\otimes 1))$. The local marginal never sees the remote POVM.
\end{proof}

\begin{theorem}[Existence of the history measure: Kolmogorov extension]
  \label{thm:exists-islimit}
  \lean{QIQTH.KolmogorovFiniteFiber.exists_isLimit}
  \leanok
  For a family of finite outcome spaces $(\alpha_i)_i$ and a Kolmogorov-consistent
  family $F$ of finite marginals, there is a probability measure $\mu$ on
  $\prod_i \alpha_i$ that is the projective limit of $F$.
\end{theorem}

\begin{proof}
  \leanok
  Each finite marginal $F.\mu(J)$ is a probability measure (\texttt{F.isProb}), hence
  finite. The candidate is $\mu = \texttt{kolmogorovMeasure}(F.\mathrm{proj})$,
  obtained by Carath\'eodory extension of the projective-family content; the work is
  to show that content is a premeasure, i.e.\ $\sigma$-subadditive.

  The crux is finite-fiber compactness. With each $\alpha_i$ carrying the discrete
  topology and finite, every measurable cylinder is closed
  (\texttt{isClosed\_of\_mem\_measurableCylinders}) and, the fibers being finite,
  compact. Hence an antitone sequence of \emph{nonempty} closed cylinders cannot have
  empty intersection (the finite-intersection property,
  \texttt{exists\_eq\_empty\_of\_antitone\_isClosed\_iInter\_empty}). This yields
  continuity of the content at $\varnothing$
  (\texttt{projectiveFamilyContent\_tendsto\_zero}), which upgrades finite additivity
  to $\sigma$-subadditivity
  (\texttt{isSigmaSubadditive\_projectiveFamilyContent}). Carath\'eodory then produces
  a genuine measure, shown to restrict back to every finite marginal --- i.e.\ to be
  the projective limit --- and to be a probability measure
  (\texttt{kolmogorovMeasure\_isProjectiveLimit}).

  This is the correlated case (general consistent family), not merely the i.i.d.\
  product.
\end{proof}

\begin{theorem}[State-agnostic typicality measure]
  \label{thm:exists-typicality}
  \lean{QIQTH.StateNetMeasure.EffectStateNet.exists_typicalityMeasure}
  \leanok
  \uses{thm:exists-islimit}
  For an \emph{effect--state net} $S$ --- a positive, normalized, additive functional
  $\omega$ with effects coarse-graining consistently across a directed system of
  contexts --- there is a probability measure $\mu$ on the global history space
  realizing all of $S$'s finite Born marginals as a projective limit. Only a normal
  state is needed, not a particular density matrix.
\end{theorem}

\begin{proof}
  \leanok
  The net's finite Born marginals assemble into a Kolmogorov-consistent family
  $S.\texttt{toFiniteMarginals}$ (positivity and coarse-graining give exactly the
  consistency the projective family requires). Applying
  Theorem~\ref{thm:exists-islimit} to it yields $\mu$ directly.
\end{proof}


\chapter{Phase B: Normal States on $\BH$ and a Free-Field Instance}
\label{chap:phaseB}

Phase~B discharges the abstract ``normal state'' input with a concrete construction
on $\BH$ for separable $\HH$, closes the typicality loop end-to-end on $\BH$, and
exhibits a relativistic free-field instance with a genuine boost symmetry.

\begin{definition}[Diagonal normal state]
  \label{def:diag-state}
  For an orthonormal basis $(b_i)$ of $\HH$ and a summable weight $(p_i)$, set
  $\omega(x) = \sum_i p_i\,\re\langle b_i, x\,b_i\rangle$ for $x \in \BH$.
\end{definition}

\begin{theorem}[The diagonal state is a normal state]
  \label{thm:normal-state}
  \lean{QIQTH.NormalState.diagStateHom_one}
  \leanok
  \uses{def:diag-state}
  If $(b_i)$ is orthonormal, $p_i \ge 0$, $(p_i)$ summable and $\sum_i p_i = 1$, then
  $\omega$ is a positive, additive functional on $\BH$ with $\omega(1) = 1$ --- a
  genuine normal state.
\end{theorem}

\begin{proof}
  \leanok
  Additivity and the bound making the defining series summable are packaged in
  \texttt{diagStateHom}; only normalization remains. Evaluate at the identity: for the
  unit vector $b_i$,
  \[
    \re\langle b_i, 1\cdot b_i\rangle = \re\langle b_i, b_i\rangle = \|b_i\|^2 = 1
  \]
  (\texttt{ContinuousLinearMap.one\_apply}, \texttt{inner\_self\_eq\_norm\_sq}, and
  orthonormality \texttt{hb.norm\_eq\_one}). Hence termwise $p_i \cdot 1 = p_i$, so
  $\omega(1) = \sum_i p_i$ (\texttt{diagState\_one}); with $\sum_i p_i = 1$ this is
  $\omega(1) = 1$.
\end{proof}

\begin{theorem}[Typicality measure on $\BH$, end to end]
  \label{thm:bh-typicality}
  \lean{QIQTH.BHTypicalityMeasure.bh_typicalityMeasure_exists}
  \leanok
  \uses{thm:normal-state, thm:exists-typicality}
  For a diagonal normal state on $\BH$ and finite, discrete record spaces
  $(\alpha_i)$, there is a probability measure $\mu$ on $\prod_i \alpha_i$ realizing
  the state's Born marginals as a projective limit.
\end{theorem}

\begin{proof}
  \leanok
  The diagonal normal state (Theorem~\ref{thm:normal-state}), paired with the
  deterministic record effects, assembles into an effect--state net
  \texttt{diagNet} on $\BH$. Its typicality measure exists by
  Theorem~\ref{thm:exists-typicality} applied to \texttt{diagNet}. The Layer-C
  construction is thereby instantiated on a genuine operator-algebraic state.
\end{proof}

\begin{theorem}[Boost invariance of the free-field measure]
  \label{thm:freefield-boost}
  \lean{QIQTH.FreeFieldTypicality.freeFieldMeasure_boost_invariant}
  \leanok
  \uses{thm:bh-typicality}
  Model a free field by occupation sectors $m \to \mathrm{Bool}$ over modes $m$, with a
  boost acting as a permutation $e: m \simeq m$ via
  $(\mathrm{boost}_e\, s)(i) = s(e\,i)$. If the single-mode law $\nu$ is
  boost-invariant, then so is the global history measure:
  $(\mu_\infty)\circ \mathrm{boost}_e^{-1} = \mu_\infty$.
\end{theorem}

\begin{proof}
  \leanok
  The global measure is the countable product $\mu_\infty = \bigotimes_{\mathbb{N}}
  \nu$ (\texttt{Measure.infinitePi}). By the projective-limit uniqueness criterion
  (\texttt{isLimit\_map\_eq}) it suffices to match, for every finite coordinate set
  $J$, the $J$-marginal of the pushforward with the finite product $\prod_J \nu$.

  Restricting to $J$ after a global boost equals applying the boost coordinatewise
  after restricting: $J\text{-restrict} \circ \mathrm{boost}_e =
  (\text{coordinatewise } \mathrm{boost}_e) \circ J\text{-restrict}$. Pushing forward
  along this composite (\texttt{Measure.map\_map}, with measurability of the boost and
  of the restriction) turns the $J$-marginal of $\mu_\infty$, which is $\prod_J \nu$
  (\texttt{freeFieldMeasure\_marginal}), into the product of the pushed-forward factors
  (\texttt{Measure.pi\_map\_pi}). Each factor is invariant by the hypothesis
  $\nu\circ\mathrm{boost}_e^{-1} = \nu$ (which also supplies the $\sigma$-finiteness
  side conditions). So the $J$-marginal is again $\prod_J \nu$, and invariance
  follows.
\end{proof}


\chapter{Scope and Open Frontiers}

\noindent\textbf{What is \emph{not} claimed.}
\begin{itemize}
  \item \emph{Unconditional Born.} Layer~B proves Born as a conditional
        representation/typicality theorem and proves its premises necessary; it does
        not derive probability from nothing.
  \item \emph{The continuum (type~III$_1$) prize.} The measure/state/covariance
        machinery above is complete and runs on a finite-mode free-field net with a
        real boost symmetry; lifting it to the continuum needs Fock/CCR/quasifree
        infrastructure not yet in Mathlib. That specific local field algebras are
        type~III$_1$ (Buchholz--Wichmann) is \emph{cited}, not proved.
\end{itemize}
